\documentclass{reportkit}
\usepackage{float}

\setreportkitleftheader{PHÂN TÍCH NGHỀ NGHIỆP CÁ NHÂN}
\setreportkitfooter{Con đường vào AI \& Data Engineering tại Úc}
\setreportkitversion{v1.0}

\title{Kim Chỉ Nam Nghề Nghiệp: Con Đường Vào AI \& Data Engineering Tại Úc}
\author{Áp dụng khung CR = D × C × E × A × R vào hồ sơ cá nhân}
\date{4 tháng 9, 2026}

\begin{document}
\maketitle

\begin{execsummary}
Hồ sơ hiện tại có nền tảng thực chiến AI hiếm gặp ở sinh viên/tân cử nhân, đặc biệt là
kỳ thực tập tại VinBigdata với bằng chứng ở tầng ứng dụng thực tế (fine-tuning LLM,
ASR, NLU). Vấn đề không nằm ở năng lực chuyên môn, mà ở hai điểm: (1) bằng chứng --
portfolio GitHub hiện dừng ở mức phân tích độc lập, chưa phản ánh đúng tầm của kinh
nghiệm đã có; và (2) khoảng trống kỹ năng triển khai/vận hành (MLOps), yếu tố quyết
định mức lương cao nhất tại thị trường Úc. Định hướng đề xuất là AI Engineering /
MLOps thay vì ML Research thuần túy. Ưu tiên 90 ngày tới: nâng cấp một dự án hiện có
lên chuẩn triển khai được, thay vì học thêm một khóa học hay làm thêm một dự án mới.
\end{execsummary}

\section{Giả thuyết nghề nghiệp}
Áp dụng khung Professional Career Exploration Framework, giả thuyết nghề nghiệp hiện
tại được phát biểu như sau, cùng với hai câu hỏi cần kiểm chứng trước khi cam kết
hoàn toàn:

\begin{deliverablenote}{Giả thuyết nghề nghiệp (đã điền)}
``Tôi muốn trở thành một AI/ML Engineer vì đây là ngành có nhu cầu tuyển dụng tăng
nhanh nhất tại Úc và đã có nền tảng thực chiến từ kỳ thực tập AI. Công việc này đòi
hỏi không chỉ xây dựng mô hình mà còn triển khai và vận hành hệ thống AI trong môi
trường thực tế. Trước khi cam kết hoàn toàn, cần kiểm chứng: liệu hướng đi phù hợp là
kỹ sư triển khai (MLOps/AI Engineering) hay hướng nghiên cứu mô hình (ML Research)?''
\end{deliverablenote}

Dựa trên bằng chứng hiện có, hướng \textbf{triển khai/vận hành (MLOps, AI
Engineering)} phù hợp hơn với hồ sơ hiện tại -- đây cũng là hướng có mức lương cao
nhất và ít cạnh tranh hơn hướng nghiên cứu thuần túy, vốn thường yêu cầu bằng Tiến sĩ.

\section{Bức tranh năng lực hiện có}
Điểm mạnh nổi bật nhất trong hồ sơ là kỳ thực tập AI Intern tại VinBigdata -- bằng
chứng ở tầng ``Application'' (ứng dụng thực tế), hiếm gặp ở sinh viên/tân cử nhân
(Bảng~\ref{tab:evidence}).

\begin{table}[H]
\caption{Năng lực hiện có và bằng chứng cụ thể}
\label{tab:evidence}
\centering
\begin{tabular}{p{3.1cm}p{7.3cm}l}
\toprule
Năng lực & Bằng chứng cụ thể & Tầng bằng chứng \\
\midrule
Fine-tuning mô hình LLM & Tối ưu ViGPT (1.6B--13B tham số) cho Q\&A/dialogue, giảm 15\%
  thời gian suy luận & Application \\
ASR / xử lý giọng nói & Gán nhãn 30K+ giờ dữ liệu, nâng độ chính xác nhận dạng lên
  98\% & Application \\
NLU / phát hiện ý định & Tinh chỉnh mô hình, cải thiện độ chính xác 10\% trong kịch
  bản lái xe thực tế & Application \\
Xử lý dữ liệu quy mô lớn & Xử lý 600GB+ dữ liệu văn bản tiếng Việt & Application \\
SQL, dự báo (ARIMA, XGBoost) & Dự báo nhu cầu hỗ trợ lập kế hoạch thu mua tại
  Matthews Liquor & Application \\
Tự động hóa / công cụ agentic & Xây dựng báo cáo tự động bằng Claude Code, tiết kiệm
  10 giờ thủ công/tuần & Application \\
Giao tiếp đa bên liên quan & Làm việc với đội ngũ lâm sàng/QA (Twin Global), 13 cửa
  hàng (Matthews Liquor) & Application \\
\bottomrule
\end{tabular}
\source{Trích từ hồ sơ cá nhân; phân loại tầng bằng chứng theo khung đánh giá năng lực.}
\end{table}

\begin{evidencenote}{Tổ hợp năng lực thực chất}
Đây là tổ hợp năng lực thực chất hơn nhiều so với đa số hồ sơ sinh viên cùng trình độ.
Vấn đề không nằm ở việc ``chưa đủ giỏi'', mà ở việc hồ sơ minh chứng (portfolio) chưa
phản ánh đúng tầm của những kinh nghiệm này.
\end{evidencenote}

\section{Khoảng trống năng lực}
So với mô hình năng lực thị trường Úc -- nơi trả lương cao nhất cho khả năng triển
khai và vận hành mô hình, không chỉ xây dựng -- ba khoảng trống chính nổi lên
(Bảng~\ref{tab:gap}).

\begin{table}[H]
\caption{Khoảng trống năng lực so với thị trường}
\label{tab:gap}
\centering
\begin{tabular}{p{3.6cm}p{2.6cm}p{6.1cm}}
\toprule
Năng lực còn thiếu & Mức độ quan trọng & Ghi chú \\
\midrule
Cloud (AWS/GCP/Azure) & Cao & Chưa có bằng chứng trong hồ sơ \\
MLOps / triển khai (Docker, CI/CD, giám sát mô hình) & \textbf{Cao nhất} & Đây là yếu
  tố quyết định mức lương -- chưa có bằng chứng, đây là nút thắt chính \\
RAG / vector database / agent framework & Đang tăng mạnh & Chưa có, nhưng gần với
  kinh nghiệm LLM đã có \\
\bottomrule
\end{tabular}
\source{So sánh với yêu cầu thị trường AI/ML Engineer tại Úc.}
\end{table}

Để chẩn đoán có hệ thống, áp dụng lại phương trình mức độ sẵn sàng nghề nghiệp từ
khung gốc:

\begin{equation}
  \mathrm{CR} = D \times C \times E \times A \times R,
  \label{eq:readiness}
\end{equation}

trong đó $D$ là định hướng, $C$ là năng lực nền, $E$ là bằng chứng, $A$ là khả năng
tiếp cận thị trường tuyển dụng (bao gồm visa), và $R$ là kỹ năng tuyển dụng (hồ sơ ứng
tuyển và phỏng vấn).

\begin{researchproblem}{Yếu tố nào đang gần bằng 0?}
Theo Phương trình~\ref{eq:readiness}, định hướng ($D$) và năng lực nền ($C$) không
phải điểm yếu trong hồ sơ này. Yếu tố gần bằng 0 nhất chính là \textbf{bằng chứng
($E$)} cho khả năng triển khai sản phẩm AI vào môi trường thực tế. Đây là yếu tố cần
ưu tiên xử lý tiếp theo -- không phải học thêm một khóa học ML nữa.
\end{researchproblem}

\section{Đánh giá portfolio GitHub hiện tại}
Ba dự án hiện tại (ví dụ: Customer Churn Analysis) đang dừng ở tầng ``Independent
Analysis'' -- phân tích tốt, có framing kinh doanh rõ ràng, so sánh nhiều mô hình (đạt
86\% accuracy, AUC-ROC 0.84) -- nhưng chưa ở mức ``professionally realistic''
(Hình~\ref{fig:portfolio}).

\begin{diagram}[
  type=network,
  caption={Vị trí hiện tại của portfolio so với mục tiêu cần đạt.},
  label={fig:portfolio},
  source={Minh họa khái niệm dựa trên thang mức độ thực tế của dự án.},
  description={Năm mức độ theo thứ tự, vị trí hiện tại được đánh dấu ở mức Independent, mục tiêu ở mức Professionally realistic.}
]
  \RKNode[text width=17mm]{guided}{0.2}{0}{Guided}
  \RKNode[text width=17mm]{modified}{3.1}{0}{Modified}
  \RKNode[rk accent,text width=19mm]{independent}{6.2}{0}{Independent\\(hiện tại)}
  \RKNode[rk accent,text width=19mm]{realistic}{9.5}{0}{Realistic\\(mục tiêu)}
  \RKNode[text width=17mm]{validated}{12.6}{0}{Validated}
  \RKEdge{sequence}{guided}{modified}
  \RKEdge{sequence}{modified}{independent}
  \RKEdge{sequence}{independent}{realistic}
  \RKEdge{sequence}{realistic}{validated}
\end{diagram}

\begin{limitationnote}{Khoảng cách cụ thể cần đóng}
Toàn bộ nằm trong 1 file Jupyter notebook, không có cấu trúc \code{src/}; không có
\code{requirements.txt} nên người khác không thể tái tạo môi trường; không có kiểm
thử (test) hay theo dõi thực nghiệm (experiment tracking); và không có bất kỳ hình
thức triển khai nào (API, container, pipeline).
\end{limitationnote}

\begin{principle}{Đào sâu một dự án, không dàn trải}
Ba dự án cùng một tầng bằng chứng không cộng dồn thành một dự án sâu hơn. Thay vì làm
thêm dự án thứ tư, nên nâng cấp \emph{một} dự án hiện có lên tầng kế tiếp.
\end{principle}

\section{Vai trò mục tiêu tại thị trường Úc}
Bối cảnh thị trường: AI Engineer hiện là nghề tăng trưởng nhanh nhất tại Úc theo
LinkedIn (tăng khoảng 150\%), với mức thiếu hụt nhân lực AI dự kiến lên tới 60.000
người vào năm 2027. Nhu cầu rất lớn, nhưng phần thưởng cao nhất dành cho người có khả
năng đưa mô hình vào vận hành thực tế (Bảng~\ref{tab:roles}).

\begin{table}[H]
\caption{Đánh giá mức độ phù hợp theo vai trò mục tiêu}
\label{tab:roles}
\centering
\begin{tabular}{p{3.4cm}p{5.6cm}p{3.3cm}}
\toprule
Vai trò & Mức độ phù hợp & Mức lương tham khảo (Úc) \\
\midrule
ML Engineer (junior/grad) & Phù hợp trực tiếp -- kinh nghiệm VinBigdata hỗ trợ mạnh &
  \textasciitilde148K--170K+ AUD (mid-level) \\
Data/AI Engineer & Phù hợp -- kinh nghiệm pipeline, SQL, dự báo hỗ trợ tốt & Tương tự,
  phổ rộng cơ hội hơn \\
AI Engineer (yêu cầu MLOps rõ ràng) & Chưa nên nhắm tới ngay -- cần đóng khoảng trống
  triển khai trước & Tránh bị loại ở vòng phỏng vấn kỹ thuật \\
\bottomrule
\end{tabular}
\source{Ước tính dựa trên khảo sát thị trường lao động AI/ML tại Úc; mang tính minh họa.}
\end{table}

\section{Yếu tố khả năng tiếp cận -- visa và di trú}
Với vai trò là du học sinh quốc tế, yếu tố tiếp cận ($A$ trong Phương
trình~\ref{eq:readiness}) thường là nút thắt thực sự, độc lập với năng lực chuyên môn
(Hình~\ref{fig:visa}).

\begin{diagram}[
  type=network,
  caption={Trình tự các mốc visa cần xác nhận sớm.},
  label={fig:visa},
  source={Minh họa khái niệm; không thay thế tư vấn di trú chính thức.},
  description={Bốn mốc theo trình tự: kết thúc khóa học, nộp visa 485, đánh giá kỹ năng ACS, phương án dự phòng.}
]
  \RKNode[text width=22mm]{finish}{0.2}{0}{Kết thúc\\khóa học}
  \RKNode[text width=22mm]{visa485}{4.0}{0}{Nộp visa 485\\(trong 6 tháng)}
  \RKNode[text width=22mm]{acs}{7.8}{0}{Đánh giá\\kỹ năng ACS}
  \RKNode[rk accent,text width=22mm]{fallback}{11.6}{0}{Phương án\\dự phòng}
  \RKEdge{sequence}{finish}{visa485}
  \RKEdge{sequence}{visa485}{acs}
  \RKEdge{sequence}{acs}{fallback}
\end{diagram}

\begin{decisionpoint}{Lộ trình visa nào là phương án chính?}
\textbf{Visa 485 (Temporary Graduate)}: cầu nối chuẩn sau khi tốt nghiệp -- cần nộp
trong vòng 6 tháng sau khi hoàn thành khóa học, đang ở Úc tại thời điểm nộp.
\textbf{Đánh giá kỹ năng ACS}: bắt buộc cho các visa định cư diện kỹ năng
(189/190/491, 494/186) -- giá trị 2 năm, cần chọn đúng mã ngành nghề ANZSCO ngay từ
đầu vì chọn sai có thể bị đánh giá thấp hơn thực tế. \textbf{Bảo lãnh từ doanh nghiệp
(482/186)}: phương án khả thi hơn sau ``vách đá 485'' nếu các visa tính điểm quá cạnh
tranh. Nên xác nhận sớm mốc thời gian cho cả bốn bước ở Hình~\ref{fig:visa}.
\end{decisionpoint}

\section{Kế hoạch hành động ưu tiên (90 ngày)}
\begin{enumerate}
  \item \textbf{Nâng cấp một dự án hiện có} (ví dụ: Customer Churn Analysis) thành sản
    phẩm triển khai được:
    \begin{itemize}
      \item Tách notebook thành cấu trúc \code{src/} rõ ràng
      \item Thêm \code{requirements.txt}, lưu model artifact
      \item Bọc inference bằng API đơn giản (FastAPI)
      \item Đóng gói bằng Docker, triển khai miễn phí (Render/Railway/AWS free tier)
      \item Viết README chuẩn, thêm test cơ bản
    \end{itemize}
  \item \textbf{Kết hợp mảng mạnh nhất (LLM/NLU từ VinBigdata) với kỹ năng triển khai
    còn thiếu}: xây một dự án RAG hoặc agent nhỏ, deploy lên cloud, có giám sát cơ bản
    -- đây là dự án bắc cầu giữa kinh nghiệm cũ và yêu cầu thị trường Úc hiện tại.
  \item \textbf{Xác nhận lộ trình visa} song song với kế hoạch kỹ thuật -- không để
    yếu tố này trở thành nút thắt bị phát hiện quá muộn.
  \item \textbf{Thu thập 30--50 tin tuyển dụng AI/ML Engineer tại Úc}, gắn nhãn theo:
    yêu cầu kỹ thuật, có bảo lãnh visa hay không, cấp độ kinh nghiệm -- dùng làm cơ sở
    nhắm mục tiêu ứng tuyển thực tế thay vì phỏng đoán.
\end{enumerate}

\section{Vòng lặp phản hồi liên tục}
Tài liệu này nên được xem lại định kỳ khi có dự án mới, phản hồi phỏng vấn, hoặc thay
đổi về tình trạng visa -- đúng theo tinh thần của khung ``vòng lặp phản hồi liên tục''
gốc (Hình~\ref{fig:loop}).

\begin{diagram}[
  type=cycle,
  caption={Vòng lặp cập nhật hồ sơ nghề nghiệp.},
  label={fig:loop},
  source={Minh họa khái niệm.},
  description={Bốn giai đoạn theo vòng tròn: hành động, thu thập phản hồi, cập nhật bằng chứng, đánh giá lại giả thuyết, quay lại hành động.}
]
  \RKCycleSetup{3.8}
  \RKCycleNode{action}{90}{Hành động\\(dự án, ứng tuyển)}
  \RKCycleNode[rk accent]{feedback}{-18}{Thu thập\\phản hồi}
  \RKCycleNode{evidence}{-162}{Cập nhật\\bằng chứng}
  \RKCycleEdge{action}{feedback}
  \RKCycleEdge{feedback}{evidence}
  \RKCycleEdge[quay lại]{evidence}{action}
\end{diagram}

\end{document}
